\documentclass[11pt]{article}
\usepackage{geometry}
\geometry{paperwidth=4.25in,paperheight=5in,margin=0.25in}
\usepackage{tikz}

%\usepackage[notextcomp]{kpfonts}
%\RequirePackage[T1]{fontenc}

\newcommand{\drawfoot}{%
  \begin{scope}[xshift=0.06in,color=black]
    \draw[fill=black!70!white] (0,0) arc (0:180:0.12in) -- (0,0) -- cycle;
    \draw (-0.1in,0) arc (0:30:0.12in);
  \end{scope}
}

\newcommand{\drawhorn}{%
      \draw[color=brown!80!black,fill=brown] (0,0)
      arc (30:60:0.45in)
      arc (30:0:0.4in);
}

\newcommand{\draweye}{%
      \draw[fill=white] (0,0) circle (0.09in);
      \draw[fill=black] (0.02in,-0.02in) circle (0.02in);
}

\newcommand{\headpath}{%
(0,0in)
arc (90:175:0.35in) % top ear
arc (250:300:0.25in) % bottom ear
arc (170:270:0.30in) % chin
arc (-90:10:0.30in) % chin
arc (-120:-70:0.25in) % bottom ear
arc (5:90:0.35in) % top ear
-- cycle
}

\setlength{\parindent}{0in}

\begin{document}
\sffamily

\begin{center}
  \begin{tikzpicture}[line width=1pt,scale=0.6, every node/.style={transform shape}]
    % draw body
    \begin{scope}
      \draw[line width=1pt] (0,0) circle (0.5in);
      
      \clip (0,0) circle (0.5in);
      \begin{scope}[color=black!85!white]
        
    \fill[xshift=0.2in,yshift=0.1in] plot [smooth cycle, tension=2] coordinates { (0,0) (0.6,0.6) (-0.6,0.6) };

    \fill[xshift=0.4in,yshift=-0.4in] plot [smooth cycle, tension=2] coordinates { (0,0) (0.6,0.6) (-0.6,0.6) };

    \fill[xshift=-0.4in,yshift=-0.6in] plot [smooth cycle, tension=2] coordinates { (0,0) (0.6,0.6) (-0.6,0.6) };

    \fill[xshift=-0.2in,yshift=0.1in] plot [smooth cycle, tension=2] coordinates { (0,0) (-0.2,-0.6) (0.6,-0.6) };

    \fill[xshift=-0.2in,yshift=0.5in] plot [smooth cycle, tension=2] coordinates { (0.35,0.35) (-0.9,-0.8) (0.2,-0.2) };

      \end{scope}

    \end{scope}

    \begin{scope}
      \clip (0,0) circle (0.5in);

      \fill[color=black,join=round] (1.75,-1.75) -- (1.75,0.5) --
        (3.75, 0.5) -- (3.75,-1.75);
    \end{scope}

    % draw right foot
    \begin{scope}[yshift=-0.53in]
      \begin{scope}[xshift=0.3in]
        \drawfoot
      \end{scope}

      % draw left foot
      \begin{scope}[xshift=-0.3in,xscale=-1]
        \drawfoot
      \end{scope}
    \end{scope}

    % draw head
    \begin{scope}[yshift=0.6in,xshift=-.15in]
    \draw[fill=white] \headpath;

    \clip \headpath;
    \fill[color=black!85!white,xshift=0.65in,yshift=0.2in] plot [smooth cycle, tension=2] coordinates { (0,0) (-1,-1) (-0.5,0.5) };

    % draw nose
    \begin{scope}[xshift=.15in,xscale=1.5,yshift=-0.6in,color=pink!50!black]
      \draw[fill=pink] (0in,0) circle (0.09in);
      \draw (0.1,0.1) -- (0.1,0);
      \draw (-0.1,0.1) -- (-0.1,0);
    \end{scope}

    \draw[stroke=black,fill=none] \headpath;
    \end{scope}

    % draw horn
    \begin{scope}[yshift=0.55in]
      \begin{scope}[xshift=-.15in]
        \drawhorn
      \end{scope}
      \begin{scope}[xshift=.15in,xscale=-1]
        \drawhorn
      \end{scope}
    \end{scope}

    % draw eyes
    \begin{scope}[yshift=0.4in]
      \begin{scope}[xshift=-.11in]
        \draweye
      \end{scope}
      \begin{scope}[xshift=.11in,xscale=-1]
        \draweye
      \end{scope}
    \end{scope}

    % draw title MOOCulus
    \node[anchor=west] at (0.5in,0.2) {\scalebox{3.2}{\sffamily\Huge\textbf{M}}};
    \node[anchor=west] at (1.5in,0.2) {\scalebox{3.2}{\sffamily\Huge\textbf{O}}};
    \node[anchor=west] at (2.3in,0.2) {\scalebox{3.2}{\sffamily\Huge\textbf{O}}};
    \node[anchor=west] at (3.1in,0.2) {\scalebox{3.2}{\sffamily\Huge\textbf{C}}};
    \node[anchor=west] at (3.8in,0.2) {\scalebox{3.2}{\sffamily\Huge ulus}};

    % draw subtitle
    \node[anchor=north west] at (0.6in,-0.30in) {\sffamily\scalebox{2.1}{\footnotesize massive open online calculus}};

    
  \end{tikzpicture}
\end{center}

\vspace{-3ex}

\begin{tabular}{@{}c@{}c@{}}%
\begin{minipage}{0.5\textwidth}%
\begin{center}
\includegraphics[width=\textwidth]{mooculus-qrcode.pdf}

\vspace{-4ex}
\scalebox{0.8}{\textsf{https://mooculus.osu.edu/}}%
\end{center}
\end{minipage} &%
\begin{minipage}{0.5\textwidth}%
\begin{center}
\huge
Free Textbook! \\
Free Videos! \\
Free Homework! \\
\end{center}
\end{minipage} \\
\end{tabular}

\vfill

Calculus is about the very large, the very small, and how things
change---the surprise is that something seemingly so abstract ends up
explaining the real world.

\vfill

This course is a first and friendly introduction to calculus, suitable
for someone who has never seen the subject before, or for someone who
has seen some calculus but wants to review the concepts and practice
applying those concepts to solve problems. One learns calculus by
doing calculus, and so \textbf{this course is based around doing
  practice problems.}



\end{document}

%%% Local Variables: 
%%% mode: latex
%%% TeX-master: t
%%% End: 
